\documentclass{fuci-article}
\usepackage[utf8]{inputenc}
\usepackage[T1]{fontenc}
\usepackage{geometry}
\usepackage{graphicx}
\usepackage{float}
\usepackage{hyperref}
\usepackage[english,indonesian]{babel}

\title{Optimasi Pemrosesan Template LaTeX Menggunakan Pendekatan Berbasis AI}
\author{Academic Support Agent}
\date{\today}

\begin{document}
\maketitle

\begin{abstract}

\end{abstract}

\chapter{PENDAHULUAN}


\chapter{TINJAUAN PUSTAKA}


% Original content (commented out):
% \documentclass{fuci-article}

\title{Mozione di indirizzo}

% Your city:
\date{???, \today}

\addbibresource{motion.bib}

\begin{document}

\makecenteredtitle

% Your group:
Il gruppo \textbf{F.U.C.I. ???} propone all'Assemblea Federale la seguente
mozione di indirizzo:

\begin{center}
  % Your motion's title here:
  \textbf{``???''}
\end{center}

\aspect{Introduzione alla mozione}

\aspect{Aspetto sociopolitico-culturale}

\aspect{Aspetto teologico-culturale}

\aspect{Aspetto universitario}

\aspect{Bibliografia}

\nocite{*}
\printbibliography[heading=none]

\vspace{2pt}

\begin{flushright}
  Per il gruppo, \\
  % Your full name
\end{flushright}

\pagebreak

I membri del gruppo

\signaturelist{39}

\end{document}


\end{document}